\documentclass{article}
\usepackage[T1]{fontenc}
%\usepackage{gensymb}
%\usepackage{microtype}
\usepackage[margin=3cm]{geometry}
\usepackage{lmodern}
\usepackage{amsmath}
\usepackage{titlesec}
\usepackage{enumitem}
% {left-indent}{space before}{space after}
\titlespacing*{\subsection}{0pt}{2.5em plus 1ex minus .2ex}{2ex plus .2ex}
\titlespacing*{\section}{0pt}{6ex plus 1ex minus .2ex}{2.3ex plus .2ex}
\linespread{1.2}
\begin{document}

\section{Log Questions}
\begin{enumerate}[label=\textbf{\arabic*.}, itemsep=1.2em]
  \item $\log_{3}(81) = 4$
  \item $\log_{5}(\frac{1}{25}) = -2$ 
  \item solve for x: \\ $\log_{2}(x) = 5 \\ x = 32$
  \item solve for b: $\log_{b}(8) = 3 \\ b = 2$
  \item $\log_{b}(1) = 0$ for every base b because $b^0 = 1$ 
  \item At a glance $\log_{2}(9)$ can't be true, because no even number can't be multiplied by itself to become odd
  \item $\log_{2}(0.5) = -1$ 
\end{enumerate}

\section{Log Rules}
\subsection{Multiplication}
If $\log_{b}(a) = m$ and $\log_{b}(c) = n$, what is the $\log_{b}(a \times c)$?
\begin{gather*}
  a = b^m \qquad \text{and} \qquad c = b^n \\
  \begin{aligned}
    a \times c &= b^m \times b^n = b^{m+n} \\
    \log_{b}(a \times c) &= m + n
  \end{aligned}
\end{gather*}

\subsection{Division}
If $log_{b}(a) = m$ and $log_{b}(c) = n$, what is the $log_{b}(a \div c)$?
\begin{gather*}
  a = b^m \qquad\text{and}\qquad c = b^n \\
  \begin{aligned}
    a \div c &= b^m \div b^n = b^{m - n} \\
    \log_{b}(a \div c) &= m - n
  \end{aligned}
\end{gather*}

\subsection{Exponentiation}
If $\log_{b}(a) = m$, what is $\log_{b}(a^k)$?
\begin{align*}
  a &= b^m \\
  \log_{b}(a^k) &= \log_{b}((a)^k) \\ &= \log_{b}(b^{km}) = km && \text{(sub m for $\log_{b}(a)$ for $a$} \\
  &= k\log_{b}(a)
\end{align*}

\section{Log Rules Questions}
\begin{enumerate}[label=\textbf{\arabic*.}, itemsep=1.2em]
  \item 
    \begin{align*}
      \text{a)} \log_{2}(8 \times 16) &= \log_{2}(128) \\
                       &= 7 \\ 
      \text{b)} \log_{2}(8 \times 16) &= \log_{2}(8) + \log_{2}(16) \\ 
                       &= 3 + 4 = 7
    \end{align*}
  \item $\log_{3}(27 \times 81) = \log_{3}(27) + \log_{3}(81) = 3 + 4 = 7$
  \item $\log_{2}(8^5) = 5\log_{2}(8) = 5 \times 3 = 15$
  \item $\log_{5}(25^3) = 3\log_{5}(25) = 3 \times 2 = 6$
  \item compress and evaluate: $\log_{2}(8) + \log_{2}(4) = \log_{2}(8 \times 4) = \log_{2}(32) = 5$
  \item $log_{b}(x^2y^3) = \log_{b}(x^2) + \log_{b}(y^3) = 2\log_{b}(x) + 3\log_{b}(y)$
  \item a) $\log_{2}(8 + 8) = \log_{2}(16) = 4$ \\ b) $\log_{2}(8) + \log_{2}(8) = 3 + 3 = 6$ \\ Therefore: $\log_{2}(8 + 8) \neq \log_{2}(8) + \log_{2}(8)$
  \item The rules for logs are derived from the laws for exponents and there are no laws for sums. 
\end{enumerate}

\end{document}
